\section{USAMO 2003}
        \begin{enumerate}
        	\item (\textit{USAMO 2003}) Prove that for every positive integer $n$ there exists an $n$-digit number divisible by $5^n$ all of whose digits are odd.


 



\textbf{Solution 1}

We proceed by induction.  For our base case, $n=1$, we have the number 5.  Now, suppose that there exists some number $a \cdot 5^{n-1}$ with $n-1$ digits, all of which are odd.  It is then sufficient to prove that there exists an odd digit $k$ such that $k\cdot 10^{n-1} + a \cdot 5^{n-1} = 5^{n-1}(k \cdot 2^{n-1} + a)$ is divisible by $5^n$.  This is equivalent to proving that there exists an odd digit $k$ such that $k \cdot 2^{n-1} + a$ is divisible by 5, which is true when $k \equiv -3^{n-1}a \pmod{5}$.  Since there is an odd digit in each of the residue classes mod 5, $k$ exists and the induction is complete.

 



\textbf{Solution 2}

First, we note that there are $5^n$ $n$ digit numbers with only odd digits.  Now, we will prove that none of these numbers have the same residue mod $5^n$, and therefore one of them must be 0 mod $5^n$.

 
Proof by contradiction:  Assume we have two distinct numbers $A_1A_2A_3...A_n$ and $B_1B_2B_3...B_n$ with only odd digits that leave the same residue mod $5^n$.  Then, subtracting the larger from the smaller would yield a new number that is a multiple of $5^n$ and has only even digits.  We could then halve all of the digits in that number to get a second multiple of $5^n$ with at most n digits that only uses the digits 0 through 4.  

 
Lemma: Every multiple of $5^n$ with n digits or less has a 5 as one of its digits.  

 All numbers of this type can be written as $k5^n$.  Then, let $k$ have $x$ factors of $2$ in it.  ($x<n$, or else our number would have more than n digits).  So, we have $k5^n=a2^x5^n=a*10^x*5^{n-x}$ for some odd a.  Now $a*10^x*5^{n-x}$ is an odd multiple of 5  ($a5^{n-x}$)  with x zeroes after it, and all multiples of 5 end in 5.  Therefore, $a*10^x*5^{n-x}$ always contains a 5 as its $(x+1)^{st}$ digit, and we have proven our lemma.  

 By our lemma, our number with only the digits 0 through 4 cannot be a multiple of $5^n$, and so we have reached a contradiction.  QED

 Note:  Not only does this prove the desired claim that there exists such a number, but it also proves that there is exactly one such number.

 



\textbf{Solution 3 (Construct digits of exemplar.)}

$5^n \pmod {10} \equiv 5$.

 $2^k 5^n \pmod {10^{k +1} } \equiv  10^k 5^{n-k} \equiv 5 \cdot 10^k$.

 For any $n$, we can construct a $a = 5^n \cdot \sum_{0 \leq i < n}{2^{x_i}}$:

 $x_0 = 1$

 $a_i = 5^n \sum_{0 \leq j \leq i}{2^{x_j}}$ 

 $x_{i} = \frac{a_{i-1}}{10^{i}} + 1 \pmod 2$. (This makes the $10^{i}$ place of $a_{i}$ odd, without changing the the smaller place digits).

 $a = a_{n-1} = 5^n \cdot  \sum_{0 \leq j < n}{2^{x_j}}$ ,  with digits in place $10^i$ odd, for $0 \leq i<n$.

 Since $\sum_{0 \leq j < n}{2^{x_j}}  < 2^n$,  $a < 5^n 2^n = 10^n$.

 By construction, the digits in all the places up through  $10^{n-1}$ are odd, and since $a<10^n$, there are no other digits. 

 In fact, if $A_m = a$ that solves the case $n=m$,  then $A_{m+1} \mod 10^ m \equiv A_m$.

 Note that in some cases (like $n=5$) , both of $x_{n-1} \in \{0,1\}$ yield distinct numbers $a-5\cdot10^{n-1}$ and $a$ with all odd digits. $a-5\cdot10^{n-1}$ has $n-1<n$ digits, and so $x_{n-1}=1$ is needed for padding. 

 <i>Alternate solutions are always welcome. If you have a different, elegant solution to this problem, please \href{https://artofproblemsolving.com/wiki/index.php?title=2003\_USAMO\_Problems/Problem\_1&action=edit}{add it} to this page.</i>

 


	\item (\textit{USAMO 2003}) A convex polygon $\mathcal{P}$ in the plane is dissected into smaller convex polygons by drawing all of its diagonals. The lengths of all sides and all diagonals of the polygon $\mathcal{P}$ are rational numbers. Prove that the lengths of all sides of all polygons in the dissection are also rational numbers.


 



\textbf{Solution}

When $\mathcal{P}$ is a triangle, the problem is trivial.  Otherwise, it is sufficient to prove that any two diagonals of the polygon cut each other into rational lengths.  Let two diagonals which intersect at a point $P$ within the polygon be $AC$ and $BD$.  Since $ABCD$ is a convex quadrilateral with sides and diagonals of rational length, we consider it in isolation.

 By the \href{https://artofproblemsolving.com/wiki/index.php?title=Law\_of\_Cosines}{Law of Cosines}, $\cos BAC = \frac{AB^2 + CA^2 - BC^2}{2AB \cdot CA}$, which is rational.  Similarly, $\cos CAD$ is rational, as well as $\cos BAD = \cos (BAC + CAD) = \cos BAC \cos CAD - \sin BAC \sin CAD$.  It follows that $\sin BAC \sin CAD$ is rational.  Since $\sin^2 CAD = 1 - \cos^2 CAD$ is rational, this means that $\frac{\sin BAC \sin CAD}{\sin^2 CAD} = \frac{\sin BAC}{\sin CAD} = \frac{\sin BAP}{\sin PAD}$ is rational.  This implies that $\frac{ AB \sin BAP}{ AD \sin PAD} = \frac{\frac{1}{2} AB \cdot AP \sin BAP}{\frac{1}{2} AP \cdot AD \sin PAD} = \frac{[BAP]}{[PAD]} = \frac{BP}{PD}$ is rational.  Define $r$ to be equal to $\frac{BP}{PD}$. We know that $\frac{BP}{PD}$ is rational; hence $r$ is rational. We also have $(1+r)PD = (1+\frac{BP}{PD} )PD = PD + BP = BD$, which is, of course, rational.  It follows that $BP$ and $PD$ both have rational length, as desired.

 
<i>Alternate solutions are always welcome. If you have a different, elegant solution to this problem, please \href{https://artofproblemsolving.com/wiki/index.php?title=2003\_USAMO\_Problems/Problem\_2&action=edit}{add it} to this page.</i>

 


	\item (\textit{USAMO 2003}) Let $n \neq 0$. For every sequence of integers


 
$A = a_0,a_1,a_2,\dots, a_n$


 


 satisfying $0 \le a_i \le i$, for $i=0,\dots,n$, define another sequence


 
$t(A)= t(a_0), t(a_1), t(a_2), \dots, t(a_n)$


 


 by setting $t(a_i)$ to be the number of terms in the sequence $A$ that precede the term $a_i$ and are different from $a_i$. Show that, starting from any sequence $A$ as above, fewer than $n$ applications of the transformation $t$ lead to a sequence $B$ such that $t(B) = B$.


 



\textbf{Solution}

Consider some sequence $C = c_0, \ldots, c_n$ as the image of $A$ after $t$ has been applied some finite number of times.

 \textbf{Lemma 1.}  If $t(c_k) = c_k = j$, then ${} c_j = \cdots = c_{j+i} = \cdots = c_k = j$ ($0 \le i \le k-j$).

 \textit{Proof.} Since the $j$ terms $c_0, \ldots, c_{j-1}$ are all less than $j$, no other terms that precede $c_{k}$ can be unequal to $c_k$.  The lemma follows.

 \textbf{Lemma 2.} If $t(c_k) = c_k = j$, then $t^2(c_k) = t(c_k)$ (where $f^2(x) = f(f(x))$).

 \textit{Proof.}  Since only $i$ terms precede the term $c_i$, we will have $t^m(c_i) \le i$, for any integer $m$.  This means that we will always have $t^m (c_0), \ldots, t^m (c_{j-1}) < j$.  This means that $c_j = t(c_j) = \cdots = c_k = t(c_k)$, and the lemma follows by iteration.

 Thus we may regard a term $c_k$ as \textit{stable} if $t(c_k) = c_k$.  We will call a sequence \textit{stable} if all of its terms are stable.

 \textbf{Lemma 3.}  If $c_k =j$ is not stable, then $t(c_k) > c_k$.

 \textit{Proof.} We have $c_0, \ldots, c_{j-1} < j$, so we always have $t(c_k) \ge c_k$.  Equality implies that $c_k$ is stable.

 We will now prove the problem by induction.  For the base case, $n=1$, we have either 0,0 or 0,1, both of which are stable.  Now, suppose that $t^{n-2}(a_0), \ldots, t^{n-2}(a_{n-1})$ are stable.  We must then have $t^{n-2}(a_n) \in \{ n-2, n-1, n\}$.  If $t^{n-2}(a_n) = n$, then it is already stable.  If $t^{n-2}(a_n) = n-1$, then either it is already stable or $t^{n-1}(a_n) = n$, which is stable.  If $t^{n-2}(a_n) = t^{n-2}(a_{n-1})$, then $t^{n-2}(a_n)$ must already be stable.  The only possibilities remaining are $t^{n-2}(a_n) = n-2, t^{n-2}(a_{n-1}) = n-1$ and $t^{n-2}(a_n) = n-2, t^{n-2}(a_{n-1}) < n-2$.  In the first case, we must have $t^{n-1}(a_n)$ equal to $n-1$ or $n$, both of which will make it stable.  In the second case, we must have $t^{n-2}(a_{n-2}) = t^{n-2}(a_{n-1}) < n-2$, giving us $t^{n-1}(a_n) = n$, which will make it stable.  Thus the induction is complete.

 
<i>Alternate solutions are always welcome. If you have a different, elegant solution to this problem, please \href{https://artofproblemsolving.com/wiki/index.php?title=2003\_USAMO\_Problems/Problem\_3&action=edit}{add it} to this page.</i>

 


 


	\item (\textit{USAMO 2003}) Let $ABC$ be a triangle. A circle passing through $A$ and $B$ intersects segments $AC$ and $BC$ at $D$ and $E$, respectively. Lines $AB$ and $DE$ intersect at $F$, while lines $BD$ and $CF$ intersect at $M$. Prove that $MF = MC$ if and only if $MB\cdot MD = MC^2$.


 



\textbf{Solution 1}

Extend segment $DM$ through $M$ to $G$ such that $FG\parallel CD$.
                    \begin{figure}[H]
                        \centering
                        \includegraphics[width=0.5\textwidth]{Solution/USAMO_I_2003_4_Solution_Figure_1.png}
                    \end{figure}
                    Then $MF = MC$ if and only if quadrilateral $CDFG$ is a parallelogram, or, $FD\parallel CG$. Hence $MC = MF$ if and only if $\angle GCD = \angle FDA$, that is, $\angle FDA + \angle CGF = 180^\circ$.

 Because quadrilateral $ABED$ is cyclic, $\angle FDA = \angle ABE$. It follows that $MC = MF$ if and only if
 \[180^\circ = \angle FDA + \angle CGF = \angle ABE + \angle CGF,\]that is, quadrilateral $CBFG$ is cyclic, which is equivalent to
 \[\angle CBM = \angle CBG = \angle CFG = \angle DCF = \angle DCM.\]Because $\angle DMC = \angle CMB$, $\angle CBM = \angle DCM$ if and only if triangles $BCM$ and $CDM$ are similar, that is
 \[\frac{CM}{BM} = \frac{DM}{CM},\]or $MB\cdot MD = MC^2$.

 



\textbf{Solution 2}

We first assume that $MB\cdot MD = MC^2$. Because $\frac{MC}{MD} = \frac{MB}{MC}$ and $\angle CMD = \angle BMC$, triangles $CMD$ and $BMC$ are similar. Consequently, $\angle MCD = \angle MBC$. This exact condition can also be visualized through power of M with respect to the circumcircle of triangle $BDC$ and getting the angle condition from the alternate segment theorem.
                    \begin{figure}[H]
                        \centering
                        \includegraphics[width=0.5\textwidth]{Solution/USAMO_I_2003_4_Solution_Figure_2.png}
                    \end{figure}
                    Because quadrilateral $ABED$ is cyclic, $\angle DAE = \angle DBE$. Hence
 \[\angle FCA = \angle MCD = \angle MBC = \angle DBE = \angle DAE = \angle CAE,\]implying that $AE\parallel CF$, so $\angle AEF = \angle CFE$. Because quadrilateral $ABED$ is cyclic, $\angle ABD = \angle AED$. Hence
 \[\angle FBM = \angle ABD = \angle AED = \angle AEF = \angle CFE = \angle MFD.\]Because $\angle FBM = \angle DFM$ and $\angle FMB = \angle DMF$, triangles $BFM$ and $FDM$ are similar. Consequently, $\frac{FM}{DM} = \frac{BM}{FM}$, or $FM^2 = BM\cdot DM = CM^2$. Therefore $MC^2 = MB\cdot MD$ implies $MC = MF$.

 Now we assume that $MC = MF$. Applying Ceva's Theorem to triangle $BCF$ and cevians $BM, CA, FE$ gives
 \[\frac{BA}{AF}\cdot\frac{FM}{MC}\cdot\frac{CE}{EB} = 1,\]implying that $\frac{BA}{AF} = \frac{BE}{EC}$, so $AE\parallel CF$.

 Consequently, $\angle DCM = \angle DAE$. Because quadrilateral $ABED$ is cyclic, $\angle DAE = \angle DBE$. Hence
 \[\angle DCM = \angle DAE = \angle DBE = \angle CBM.\]Because $\angle CBM = \angle DCM$ and $\angle CMB = \angle DMC$, triangles $BCM$ and $CDM$ are similar. Consequently, $\frac{CM}{DM} = \frac{BM}{CM}$, or $CM^2 = BM\cdot DM$.

 Combining the above, we conclude that $MF = MC$ if and only if $MB\cdot MD = MC^2$.

 


 <i>Alternate solutions are always welcome. If you have a different, elegant solution to this problem, please \href{https://artofproblemsolving.com/wiki/index.php?title=2003\_USAMO\_Problems/Problem\_4&action=edit}{add it} to this page.</i>

 


	\item (\textit{USAMO 2003}) Let $a$, $b$, $c$ be positive real numbers. Prove that


 $\dfrac{(2a + b + c)^2}{2a^2 + (b + c)^2} + \dfrac{(2b + c + a)^2}{2b^2 + (c + a)^2} + \dfrac{(2c + a + b)^2}{2c^2 + (a + b)^2} \le 8.$

 



\textbf{Solution 1}

Since all terms are homogeneous, we may assume WLOG that $a + b + c = 3$.

 Then the LHS becomes $\sum \frac {(a + 3)^2}{2a^2 + (3 - a)^2} = \sum \frac {a^2 + 6a + 9}{3a^2 - 6a + 9} = \sum \left(\frac {1}{3} + \frac {8a + 6}{3a^2 - 6a + 9}\right)$.

 Notice $3a^2 - 6a + 9 = 3(a - 1)^2 + 6 \ge 6$, so $\frac {8a + 6}{3a^2 - 6a + 9} \le \frac {8a + 6}{6}$.

 So $\sum \frac {(a + 3)^2}{2a^2 + (3 - a)^2} \le \sum \left(\frac {1}{3} + \frac {8a + 6}{6}\right) = 1 + \frac {8(a + b + c) + 18}{6} = 8$, as desired.

 



\textbf{Solution 2}

Note that
 \begin{align*} (2x + y)^2 + 2(x - y)^2 &= 4x^2 + 4xy + y^2 + 2x^2 - 4xy + 2y^2 \\ &= 3(2x^2 + y^2). \end{align*}
Setting $x = a$ and $y = b + c$ yields
 \[(2a + b + c)^2 + 2(a - b - c)^2 = 3(2a^2 + (b + c)^2).\]Thus, we have
 \[\frac{(2a + b + c)^2}{2a^2 + (b + c)^2} = \frac{3(2a^2 + (b + c)^2) - 2(a - b - c)^2}{2a^2 + (b + c)^2} = 3 - \frac{2(a - b - c)^2}{2a^2 + (b + c)^2},\]and its analogous forms. Thus, the desired inequality is equivalent to
 \[\frac{(a - b - c)^2}{2a^2 + (b + c)^2} + \frac{(b - c - a)^2}{2b^2 + (c + a)^2} + \frac{(c - a - b)^2}{2c^2 + (a + b)^2}\geq\frac{1}{2}.\]Because 

 


	\item (\textit{USAMO 2003}) At the vertices of a regular hexagon are written six nonnegative integers whose sum is 2003. Bert is allowed to make moves of the following form: he may pick a vertex and replace the number written there by the absolute value of the difference between the numbers written at the two neighboring vertices. Prove that Bert can make a sequence of moves, after which the number 0 appears at all six vertices.


 



\textbf{Solution}

Assume the original numbers are $a,b,c,d,e,f$. Since $a+b+c+d+e+f$ is odd, either $a+c+e$ or $b+d+f$ must be odd. WLOG let $a+c+e$ be odd and $a\ge c\ge e \ge 0$.

 \textbf{Case 1} $a,c,e>0$. Define \textit{Operation A} as the sequence of moves from Step 1 to Step 3, shown below:

 
                    \begin{figure}[H]
                        \centering
                        \includegraphics[width=0.5\textwidth]{Solution/USAMO_I_2003_6_Solution_Figure_1.png}
                    \end{figure}
                    

 Notice that Operation A changes the numbers $a,c,e$ to $c-e,c,a-c$ and they are all nonnegative, since $a\ge c\ge e$. Their sum changes from $a+c+e$ to $a+c-e$; it decreases as long as $e\ne 0$. If we repeat Operation A enough times, its sum will decrease and eventually we will arrive at a point where at least one of the numbers in the positions originally occupied by $a,c,e$ has become a 0.

 \textbf{Case 2} $a,c>0$ and $e=0$. Define \textit{Operation B} as the sequence of moves from Step 1 to Step 3, shown below:

 
                    \begin{figure}[H]
                        \centering
                        \includegraphics[width=0.5\textwidth]{Solution/USAMO_I_2003_6_Solution_Figure_2.png}
                    \end{figure}
                    

 where in Step 3, we take the nonnegative choice of $2c-a$ or $a-2c$. $a,c,0$ is changed to either $c,2c-a,0$ or $c,a-2c,0$. If we have $c,2c-a,0$, their sum is $3c-a$ and this is less than $a+c+0$ (the original sum) unless $a=c$, but $a\ne c$ since the original sum $a+c+0$ is odd by assumption. If we have $c,a-2c,0$, their sum is $a-c$, which is less than $a+c$. Operation B applied repeatedly will cause either $a$ or $c$ to become $0$.

 \textbf{Case 3} $a>0$ and $c=e=0$. Define \textit{Operation C} as the sequence of moves from Step 1 to Step 4, shown below:

 
                    \begin{figure}[H]
                        \centering
                        \includegraphics[width=0.5\textwidth]{Solution/USAMO_I_2003_6_Solution_Figure_3.png}
                    \end{figure}
                    

 


\end{enumerate}